\section{Implementation}
\label{sec:implementation}

The implementation is a Rust workspace containing the prover, the native
verifier, the release evaluator, and a Solana SBF verifier program.  The
on-chain program uses a deliberately small production dispatcher; legacy and
diagnostic entry points are not linked into the release binary.  This section
specifies the byte formats and state transition implemented by that binary.

\subsection{Canonical public statement}
\label{sec:implementation-statement}

The atomic public statement is
\[
  x=(P,s,A,\nu,C_{\mathsf{out}},A',a,f,\Delta),
\]
where \(P\) is the pool account address, \(s\) is its sequence number, \(A\)
and \(A'\) are the current and next anchors, \(\nu\) is the nullifier,
\(C_{\mathsf{out}}\) is the output commitment, \(a\) is the asset identifier,
\(f\) is the fee, and \(\Delta\) is the deployment domain.
Table~\ref{tab:statement-wire} gives the complete 216-byte encoding.  All multibyte integers are little endian.  Each 32-byte
digest is eight canonical M31 limbs, each encoded as a four-byte integer less
than \(2^{31}-1\).  The asset identifier is also required to be canonical in
M31, and the fee is required to be less than \(2^{30}\).

\begin{table}[t]
  \centering
  \small
  \caption{Canonical atomic-payment statement payload.}
  \label{tab:statement-wire}
  \begin{tabular}{r r l}
    \toprule
    Offset & Bytes & Field \\
    \midrule
    0   & 1  & version \(=4\) \\
    1   & 1  & private Merkle depth \(=20\) \\
    2   & 6  & reserved, all zero \\
    8   & 32 & pool address \(P\) \\
    40  & 8  & sequence \(s\) \\
    48  & 32 & current anchor \(A\) \\
    80  & 32 & nullifier \(\nu\) \\
    112 & 32 & output commitment \(C_{\mathsf{out}}\) \\
    144 & 32 & output anchor \(A'\) \\
    176 & 4  & asset identifier \(a\) \\
    180 & 4  & fee \(f\) \\
    184 & 32 & deployment domain \(\Delta\) \\
    \bottomrule
  \end{tabular}
\end{table}

The transcript binding is
\[
  d_x=\operatorname{SHA256}(
  \mathtt{aspis/atomic-payment-statement/v4}\mathbin\|\mathsf{encode}(x)).
\]
On chain, \(P\) is taken from the supplied pool account and \(s\) from the
decoded pool state.  They are not accepted as caller-provided instruction
fields.  The caller-supplied \(\Delta\) must equal the pool's stored
deployment domain before verification proceeds.  This ensures that the statement checked by the proof is the statement
for the state account that the transaction will mutate.

\subsection{Proof serialization}
\label{sec:implementation-proof-wire}

\Profile{} has a fixed 6,785-byte prefix followed by five authenticated opening
sections.  The first 6,736 bytes retain the base state-only layout; \Profile{}
appends three evaluations of the zero-factor \(D\) lane and one selector byte.
Table~\ref{tab:proof-prefix} is the normative prefix map.

\begin{table}[t]
  \centering
  \small
  \caption{\Profile{} fixed proof prefix.}
  \label{tab:proof-prefix}
  \begin{tabular}{r r l}
    \toprule
    Byte range & Bytes & Contents \\
    \midrule
    0--15       & 16   & versioned proof header \\
    16--47      & 32   & mask nonce \\
    48--79      & 32   & first-phase commitment root \\
    80--111     & 32   & second-phase commitment root \\
    112--127    & 16   & initial mask claim \\
    128--4607   & 4480 & masked sumcheck \\
    4608--5951  & 1344 & 84 statement evaluations in QM31 \\
    5952--6623  & 672  & four fold-round records \\
    6624--6687  & 64   & final polynomial \\
    6688--6695  & 8    & final work nonce \\
    6696--6727  & 32   & four fold work nonces \\
    6728--6735  & 8    & batch work nonce \\
    6736--6783  & 48   & three \(D\)-lane claims in QM31 \\
    6784         & 1    & query-branch selector \\
    \bottomrule
  \end{tabular}
\end{table}

Physical position and transcript position are intentionally distinct for the
append-only extension.  The three \(D\)-lane claims are read from offsets
6736--6783 but absorbed with the other statement claims, before batch work and
the layer-combination challenge.  The selector is absorbed after the final
work nonce.  The parser checks canonical field encodings and the complete
profile header before replaying the transcript.

After the prefix, the proof contains the five sections in
Table~\ref{tab:opening-geometry}.  For a section with value width \(w\), the
wire format is
\[
  n_{16}\;\mathbin\|\;
  \bigl(\mathsf{value}_{w}\mathbin\|\mathsf{salt}_{32}\bigr)^n
  \;\mathbin\|\;m_{32}\;\mathbin\|\;
  \mathsf{frontier}_{32m},
\]
where \(n_{16}\) and \(m_{32}\) are little-endian counters.  Query indices are
derived from the transcript and are not serialized.  The parser requires the
records to match the sorted, de-duplicated derived indices, authenticates each
\(\mathsf{value}\mathbin\|\mathsf{salt}\) record against its commitment root,
and rejects any byte remaining after the fifth section.

\begin{table}[t]
  \centering
  \small
  \caption{Authenticated opening geometry.  Widths exclude the 32-byte leaf
  salt.}
  \label{tab:opening-geometry}
  \begin{tabular}{l r r}
    \toprule
    Tree & Binary depth & Value width (bytes) \\
    \midrule
    first phase, layer zero  & 17 & 416 \\
    second phase, layer zero & 17 & 192 \\
    later layer 0            & 15 & 64  \\
    later layer 1            & 13 & 64  \\
    later layer 2            & 11 & 64  \\
    \bottomrule
  \end{tabular}
\end{table}

The common transcript state is cloned into three post-final branches.  Branch
\(b\in\{0,1,2\}\) absorbs its branch label, derives an 18-query schedule, and
evaluates \(\mathsf{Good}_{\mathrm{spend}}\).  The prover and release evaluator serialize the
least-index successful branch and only that branch's openings.  The on-chain
verifier checks the serialized branch completely; least-\(\mathsf{Good}_{\mathrm{spend}}\)
selection is additionally enforced when the concrete proof is admitted by the
release evaluator.

\subsection{Accounts and proof lifecycle}
\label{sec:implementation-accounts}

Solana accounts provide persistent byte arrays with explicit ownership and
writable locking~\cite{SolanaAccounts,SolanaTransactions}.  \Profile{} uses a
proof account, a pool account, and a nullifier-marker account.

\paragraph{Proof account.}
The proof-account image is
\[
  \mathtt{ASPU}\;\mathbin\|\;
  \mathsf{proof\_len}_{32}\;\mathbin\|\;
  \mathsf{upload\_authority}_{32}\;\mathbin\|\;
  \mathsf{proof}[\mathsf{proof\_len}].
\]
The 40-byte header is followed immediately by the proof.  Tag~0 initializes a
zeroed, program-owned account; tag~1 writes an in-range chunk after checking
the upload authority.  The lifecycle deliberately permits authorized
overwrites before sealing.  It does not maintain a chunk bitmap, so the
executor compares the complete uploaded account image with the expected image
before invoking tag~62.  Tag~62 seals the account by replacing the authority
with 32 zero bytes.  Production verification requires this zero-authority
sentinel, and subsequent upload or finalization calls fail the authority
check.

For the \ProfileProofBytes-byte release proof, the 40-byte account header gives
an exact account allocation of 65,447 bytes.  The mainnet executor uses
960-byte chunks, so a fresh upload
uses \ProfileUploadTransactions{} transactions: 67 full chunks and one
127-byte final chunk.  It submits at most 16 uploads before waiting for the
window to reach finality.  Excluding program deployment, the deterministic
fresh-account schedule is therefore:
\begin{enumerate}
  \item create the 48-byte pool account and initialize it with tag~63;
  \item create the 65,447-byte proof account and initialize it with tag~0;
  \item submit \ProfileUploadTransactions{} tag~1 upload transactions; and
  \item seal the proof account with tag~62.
\end{enumerate}
Pool creation and initialization are separate transactions; proof creation
and tag~0 initialization share one transaction.  The resulting application
setup has \ProfileMainnetSetupTransactions{} transactions.  The subsequent
tag~65 verification-and-transition call is a separate transaction.  Program
deployment is also separate from this application setup.

\paragraph{Pool account.}
The 48-byte pool image is
\[
  \mathtt{ASPS}\mathbin\|1\mathbin\|0^3
  \mathbin\|s_{64}\mathbin\|A_{32}.
\]
Tag~63 accepts only a zeroed, program-owned, writable account that also signs
the initialization transaction.  It checks the anchor encoding and that
\(s+1\) does not overflow before writing the initial state.

\paragraph{Nullifier marker.}
The required marker address is the program-derived address with seeds
\[
  (\mathtt{aspis-nullifier-v1},\nu).
\]
The pool address is intentionally not a PDA seed.  It is recorded in the
72-byte marker image
\[
  \mathtt{ASPN}\mathbin\|1\mathbin\|0^3
  \mathbin\|P_{32}\mathbin\|\nu_{32}.
\]
Consequently one nullifier has one required marker address for the deployed
program, while the marker contents retain the pool to which it was applied.

\subsection{Production instruction surface}
\label{sec:implementation-dispatch}

The release binary recognizes only the lifecycle and verification tags in
Table~\ref{tab:production-tags}.  Tag~61 is retained as an explicit fixed
failure for the former diagnostic mutation path; all other tags are rejected
before account processing.  Fixed-width fields are decoded directly, without
linking the historical generic instruction dispatcher.

\begin{table}[t]
  \centering
  \small
  \caption{Minimal \Profile{} production instruction surface.}
  \label{tab:production-tags}
  \begin{tabular}{r r l}
    \toprule
    Tag & Data bytes & Operation \\
    \midrule
    0  & 5       & initialize proof account \\
    1  & \(9+n\) & upload \(n\)-byte proof chunk \\
    59 & 178     & read-only production verification \\
    60 & 137     & verify and apply, retaining proof account \\
    61 & 1       & explicit fail-closed diagnostic tag \\
    62 & 1       & finalize proof account \\
    63 & 41      & initialize pool account \\
    64 & 1       & close a finalized proof account and refund rent \\
    65 & 137     & verify, apply, and refund proof account atomically \\
    \bottomrule
  \end{tabular}
\end{table}

The 137-byte tag~65 payload consists of the tag byte, four 32-byte digests
\((A,\nu,C_{\mathsf{out}},A')\), a four-byte asset identifier, and a four-byte
fee.  Tag~60 has the same payload and a read-only proof-account ABI.  The
mainnet tag~65 transaction supplies the five accounts in
Table~\ref{tab:tag65-accounts} in exactly that order.

\begin{table}[t]
  \centering
  \small
  \caption{Tag~65 account interface.}
  \label{tab:tag65-accounts}
  \begin{tabular}{r l l}
    \toprule
    Index & Access & Account \\
    \midrule
    0 & signer, writable& finalized program-owned proof account \\
    1 & writable        & program-owned pool account \\
    2 & writable        & program-derived nullifier account \\
    3 & signer, writable& system-owned fee payer \\
    4 & read only       & executable System Program \\
    \bottomrule
  \end{tabular}
\end{table}

\subsection{Atomic tag~65 transition}
\label{sec:implementation-atomic}

The tag~65 implementation uses the following order.
\begin{enumerate}
  \item It validates ownership, signer and writable flags, account
  distinctness, the program-derived nullifier account, the current pool anchor, and
  sequence non-overflow.  A program-owned marker must be exactly 72 zero
  bytes; alternatively, the required program-derived address may be an empty System
  Program-owned account.
  \item It constructs the required statement using the actual pool address
  and decoded sequence, then computes the statement digest.
  \item It loads the finalized proof and runs the complete production
  \Profile{} verifier.  No cross-program invocation (CPI) and no account-data
  write occurs before this
  verifier returns success.
  \item If the marker account is absent, it is created at the required
  program-derived address\@.
  A prefunded empty PDA is topped up, allocated, and assigned as necessary.
  The alternative program-owned-zero path requires no System Program CPI\@.
  \item It reacquires mutable borrows, decodes and rechecks both state
  preconditions, constructs both final account images, refunds the complete
  proof-account lamport balance to the payer, writes \((P,\nu)\) to the
  marker, and replaces \((s,A)\) in the pool with \((s+1,A')\).  The proof
  account must sign, so this refund cannot consume an account whose key is not
  controlled by the transaction builder.
\end{enumerate}

The final transaction contains two Compute Budget instructions, a
1,400,000-CU limit and a 262,144-byte heap request, followed by one tag~65
application instruction.  Solana transaction rollback covers both direct
writes, the proof refund, and the marker-creation CPI, while writable locks on
the pool and required marker serialize conflicting spends
~\cite{SolanaTransactions,SolanaComputeBudget}.  A duplicate therefore
observes either the existing marker or the advanced pool state and is
rejected.  Tag~60 remains available when retaining the proof account is
required; the evaluated mainnet transition used tag~65.

The term \emph{one transaction} refers precisely to this verification and
state transition from a previously finalized proof account, including the
proof-account refund performed by tag~65.  Program deployment, pool
initialization, proof-account creation, proof upload, and proof finalization
are setup operations.  The instruction updates commitment and nullifier
state; it does not invoke a Solana Program Library (SPL) token program or
transfer application assets.
